
\subsection{Negative Energien und das Positron}
\label{subsec:negative_energien_und_positron}

Eine der überraschendsten Eigenschaften der Dirac-Gleichung%
\index{Dirac-Gleichung!negative Energielösungen}
sind ihre Lösungen mit positiver und negativer Energie.

Für ein freies Elektron gilt

\begin{equation}
	E
	=
	\pm
	\sqrt{
		p^2c^2
		+
		m_{\mathrm e}^2c^4
	}.
	\label{eq:energiezweige_dirac_positron}
\end{equation}

Die Gleichung liefert somit zwei Energiezweige:

\begin{equation}
	E_{+}
	=
	+
	\sqrt{
		p^2c^2
		+
		m_{\mathrm e}^2c^4
	}
\end{equation}

und

\begin{equation}
	E_{-}
	=
	-
	\sqrt{
		p^2c^2
		+
		m_{\mathrm e}^2c^4
	}.
\end{equation}

Die positiven Lösungen konnten unmittelbar mit gewöhnlichen
Elektronen verbunden werden. Die physikalische Bedeutung der negativen
Lösungen war dagegen zunächst völlig unklar.

\begin{PhysicsBox}[Zwei unvermeidbare Energiezweige]
	\label{box:zwei_unvermeidbare_energiezweige}
	
	Die negativen Energielösungen sind kein Rechenfehler.
	
	Sie entstehen zwangsläufig aus der relativistischen Beziehung
	
	\begin{equation}
		E^2
		=
		p^2c^2
		+
		m_{\mathrm e}^2c^4.
	\end{equation}
	
	Eine Gleichung für \(E^2\) besitzt mathematisch sowohl eine positive
	als auch eine negative Lösung für \(E\).
\end{PhysicsBox}

\subsubsection*{Das Problem der negativen Energie}

Für ein ruhendes Teilchen gilt

\begin{equation}
	p
	=
	0.
\end{equation}

Damit ergeben sich aus Gleichung
\eqref{eq:energiezweige_dirac_positron} die beiden Ruheenergien

\begin{equation}
	E
	=
	\pm m_{\mathrm e}c^2.
	\label{eq:positive_negative_ruheenergie}
\end{equation}

Der positive Wert lautet

\begin{equation}
	E_{+}
	=
	m_{\mathrm e}c^2
	\approx
	\SI{511}{\kilo\electronvolt}.
\end{equation}

Der negative Wert lautet entsprechend

\begin{equation}
	E_{-}
	=
	-m_{\mathrm e}c^2.
\end{equation}

Zwischen den beiden Energiebereichen liegt damit ein Abstand von

\begin{equation}
	\Delta E
	=
	2m_{\mathrm e}c^2
	\approx
	\SI{1.022}{\mega\electronvolt}.
	\label{eq:energieluecke_elektron_positron}
\end{equation}

Das eigentliche Problem bestand darin, dass ein Elektron scheinbar
Energie abgeben und immer weiter in Zustände mit negativer Energie
fallen könnte.

Die gewöhnliche Materie wäre dann nicht stabil. Elektronen müssten
fortlaufend Energie abstrahlen und in immer tiefere Zustände
übergehen.

Ein solches Verhalten wird jedoch nicht beobachtet.

\begin{DidacticBox}[Warum negative Energie problematisch war]
	\label{box:problem_negative_energie}
	
	Wären unbesetzte Zustände mit immer tieferer negativer Energie frei
	zugänglich, könnte ein Elektron fortlaufend Energie verlieren.
	
	Gewöhnliche Elektronen sind jedoch stabil. Sie fallen nicht spontan in
	einen unbegrenzt tiefen Energiebereich.
	
	Die negativen Lösungen verlangten deshalb nach einer neuen
	physikalischen Interpretation.
\end{DidacticBox}

\subsubsection*{Diracs historische Deutung}

Paul Dirac\index{Dirac, Paul!negative Energien} schlug zunächst vor,
dass sämtliche Zustände mit negativer Energie bereits mit Elektronen
besetzt seien.

Diese Vorstellung wurde später als Dirac-See%
\index{Dirac-See} bezeichnet.

Da Elektronen dem Pauli-Prinzip%
\index{Pauli-Prinzip!Dirac-See} unterliegen, kann ein weiteres Elektron
nicht in einen bereits besetzten Zustand übergehen. Dadurch wären die
gewöhnlichen positiven Energiezustände stabil.

Wird einem Elektron im Bereich negativer Energie genügend Energie
zugeführt, kann es in einen positiven Energiezustand angehoben werden.

Im zuvor vollständig besetzten Bereich bleibt dann eine Lücke zurück.

Diese Lücke verhält sich so, als besäße sie

\begin{itemize}
	\item die gleiche Masse wie das Elektron,
	\item eine positive elektrische Ladung und
	\item einen Spin von \(\hbar/2\).
\end{itemize}

Dirac erkannte, dass diese Lücke wie ein neues Teilchen erscheinen
musste.

\begin{HistoryBox}[Die Vorstellung des Dirac-Sees]
	\label{box:vorstellung_dirac_see}
	
	Dirac deutete die negativen Energiezustände zunächst als vollständig
	besetzten Hintergrund.
	
	Ein fehlendes Elektron in diesem Hintergrund sollte sich wie ein
	positiv geladenes Teilchen verhalten.
	
	Dieses historische Modell führte zur Vorhersage eines Antiteilchens
	des Elektrons.
\end{HistoryBox}

Die Vorstellung eines unendlich gefüllten Dirac-Sees ist heute nicht
mehr die grundlegende Beschreibung. Sie war jedoch historisch
entscheidend, weil sie den Weg zur Deutung der negativen
Energielösungen eröffnete.

\subsubsection*{Das Positron}

Das vorhergesagte Teilchen wird als Positron%
\index{Positron} bezeichnet. Es besitzt das Symbol

\begin{equation}
	e^+.
\end{equation}

Das Elektron wird entsprechend mit

\begin{equation}
	e^-
\end{equation}

bezeichnet.

Elektron und Positron besitzen dieselbe Masse:

\begin{equation}
	m_{e^+}
	=
	m_{e^-}
	=
	m_{\mathrm e}.
	\label{eq:massen_elektron_positron}
\end{equation}

Ihre elektrischen Ladungen haben denselben Betrag, aber
entgegengesetzte Vorzeichen:

\begin{equation}
	q_{e^-}
	=
	-e
\end{equation}

und

\begin{equation}
	q_{e^+}
	=
	+e.
	\label{eq:ladungen_elektron_positron}
\end{equation}

Das Positron ist damit das Antiteilchen%
\index{Antiteilchen!Positron} des Elektrons.

\begin{PhysicsBox}[Elektron und Positron]
	\label{box:elektron_und_positron}
	
	Elektron und Positron besitzen
	
	\begin{itemize}
		\item dieselbe Masse,
		\item denselben Spin und
		\item entgegengesetzte elektrische Ladungen.
	\end{itemize}
	
	Das Positron ist kein Elektron mit negativer Energie.
	
	Es ist ein reales Teilchen mit positiver Energie und positiver
	elektrischer Ladung.
\end{PhysicsBox}

Dieser letzte Punkt ist besonders wichtig. Ein beobachtetes Positron
besitzt keine negative messbare Energie.

Die negativen Lösungen der Dirac-Gleichung werden vielmehr als Hinweis
auf einen zweiten Teilchentyp interpretiert, dessen beobachtbare Energie
positiv ist.

\subsubsection*{Experimentelle Entdeckung}

Im Jahr 1932 untersuchte Carl David Anderson%
\index{Anderson, Carl David} Teilchenspuren aus der kosmischen
Strahlung\index{kosmische Strahlung}.

In einer Nebelkammer%
\index{Nebelkammer} beobachtete er die Spur eines Teilchens, das sich
in einem Magnetfeld in die entgegengesetzte Richtung zu einem Elektron
krümmte.

Aus der Krümmung der Bahn folgte, dass das Teilchen eine positive
elektrische Ladung besaß. Die Stärke der Krümmung zeigte zugleich, dass
seine Masse mit der Elektronenmasse übereinstimmte.

Damit war das Positron experimentell nachgewiesen.

\begin{HistoryBox}[Entdeckung des Positrons]
	\label{box:entdeckung_positron}
	
	Carl David Anderson entdeckte das Positron im Jahr 1932 in der
	kosmischen Strahlung.
	
	Das neue Teilchen besaß
	
	\begin{itemize}
		\item die Masse des Elektrons und
		\item eine positive elektrische Ladung.
	\end{itemize}
	
	Damit bestätigte das Experiment eine grundlegende Vorhersage der
	Dirac-Theorie.
\end{HistoryBox}

Die Entdeckung war von großer Bedeutung. Zum ersten Mal hatte eine
physikalische Gleichung überzeugend auf ein zuvor unbekanntes
Antiteilchen hingewiesen.

\subsubsection*{Die moderne Interpretation}

In der modernen Quantenfeldtheorie%
\index{Quantenfeldtheorie!Antiteilchen}
wird das Elektron nicht mehr ausschließlich als einzelnes Teilchen mit
einer festen Wellenfunktion betrachtet.

Stattdessen wird ein Elektronenfeld beschrieben. Seine Anregungen
können als Elektronen oder Positronen erscheinen.

Die negativen Energielösungen werden dabei nicht als tatsächlich
beobachtbare Elektronen mit negativer Energie verstanden.

Sie werden mathematisch so umgedeutet, dass sie Zuständen von
Antiteilchen mit

\begin{itemize}
	\item positiver Energie,
	\item entgegengesetzter Ladung und
	\item entgegengesetzten weiteren Ladungsquantenzahlen
\end{itemize}

entsprechen.

\begin{DidacticBox}[Historisches Bild und moderne Theorie]
	\label{box:historisches_bild_moderne_theorie}
	
	Der Dirac-See ist ein anschauliches historisches Modell.
	
	Die moderne Quantenfeldtheorie beschreibt Elektronen und Positronen
	dagegen als unterschiedliche Anregungen desselben quantisierten
	Feldes.
	
	Die beobachtbaren Teilchen besitzen dabei stets positive Energie.
\end{DidacticBox}

\subsubsection*{Paarbildung}

Elektron und Positron können gemeinsam erzeugt werden. Dieser Vorgang
wird als Paarbildung%
\index{Paarbildung!Elektron-Positron-Paar} bezeichnet.

Symbolisch kann er geschrieben werden als

\begin{equation}
	\gamma
	\longrightarrow
	e^-
	+
	e^+.
	\label{eq:paarbildung_ein_photon_symbolisch}
\end{equation}

Damit ein Elektron-Positron-Paar entstehen kann, muss mindestens die
gemeinsame Ruheenergie beider Teilchen bereitgestellt werden:

\begin{equation}
	E_{\gamma}
	\geq
	2m_{\mathrm e}c^2.
\end{equation}

Daraus folgt die Mindestenergie

\begin{equation}
	E_{\gamma,\mathrm{min}}
	\approx
	\SI{1.022}{\mega\electronvolt}.
	\label{eq:mindestenergie_paarbildung_positron}
\end{equation}

In der Praxis reicht die Energiebedingung allein jedoch nicht aus. Auch
der Impuls muss erhalten bleiben.

Ein einzelnes Photon kann deshalb im leeren Raum nicht ohne weiteres
in ein Elektron-Positron-Paar zerfallen.

Meist findet die Paarbildung in der Nähe eines Atomkerns statt:

\begin{equation}
	\gamma
	+
	Z
	\longrightarrow
	e^-
	+
	e^+
	+
	Z.
	\label{eq:paarbildung_atomkern}
\end{equation}

Der Atomkern \(Z\) nimmt dabei einen Teil des Impulses auf.

Alternativ können auch zwei ausreichend energiereiche Photonen
miteinander wechselwirken:

\begin{equation}
	\gamma
	+
	\gamma
	\longrightarrow
	e^-
	+
	e^+.
	\label{eq:breit_wheeler_paarbildung}
\end{equation}

\begin{PhysicsBox}[Energie wird zu Materie]
	\label{box:energie_wird_zu_materie}
	
	Bei der Paarbildung wird elektromagnetische Energie in die Masse und
	Bewegungsenergie eines Elektrons und eines Positrons umgewandelt.
	
	Die Mindestenergie entspricht der gemeinsamen Ruheenergie:
	
	\begin{equation}
		E_{\mathrm{min}}
		=
		2m_{\mathrm e}c^2
		\approx
		\SI{1.022}{\mega\electronvolt}.
	\end{equation}
	
	Energie, Impuls und elektrische Ladung bleiben dabei erhalten.
\end{PhysicsBox}
\begin{DidacticBox}[Wo entsteht ein Antiteilchen?]
	\label{box:wo_entsteht_ein_antiteilchen}
	
	Ein Positron ist nicht bereits irgendwo verborgen vorhanden und wird
	bei einem Prozess lediglich freigesetzt.
	
	Es entsteht neu, wenn genügend Energie zur Verfügung steht und alle
	Erhaltungssätze erfüllt sind. Die Energie des ursprünglichen Systems
	wird dabei unter anderem in die Ruheenergie der neu erzeugten Teilchen
	umgewandelt:
	
	\begin{equation}
		E
		\longrightarrow
		2m_{\mathrm e}c^2
		+
		E_{\mathrm{kin}}.
	\end{equation}
	
	Ein Positron kann beispielsweise entstehen
	
	\begin{itemize}
		\item bei der Paarbildung eines energiereichen Photons in der Nähe
		eines Atomkerns,
		\item bei der Wechselwirkung zweier hochenergetischer Photonen,
		\item beim \(\beta^+\)-Zerfall eines Atomkerns oder
		\item bei energiereichen Teilchenkollisionen.
	\end{itemize}
	
	In einem gewöhnlichen Wasserstoffatom befindet sich dagegen kein reales
	Positron. Das Atom enthält ein Proton und ein gebundenes Elektron. Die
	Möglichkeit von Antiteilchen gehört zur Struktur des Elektronenfeldes,
	bedeutet aber nicht, dass in jedem Atom ein Positron vorhanden sein
	muss.
\end{DidacticBox}
\subsubsection*{Annihilation}

Treffen ein Elektron und ein Positron aufeinander, können sie sich
gegenseitig vernichten.

Dieser Vorgang wird als Annihilation%
\index{Annihilation!Elektron und Positron} bezeichnet.

Im einfachsten Fall entstehen zwei Photonen:

\begin{equation}
	e^-
	+
	e^+
	\longrightarrow
	\gamma
	+
	\gamma.
	\label{eq:elektron_positron_annihilation}
\end{equation}

Befinden sich Elektron und Positron vor der Annihilation nahezu in
Ruhe, besitzt jedes Photon die Energie

\begin{equation}
	E_{\gamma}
	=
	m_{\mathrm e}c^2
	\approx
	\SI{511}{\kilo\electronvolt}.
	\label{eq:energie_annihilationsphoton}
\end{equation}

Die beiden Photonen bewegen sich dann in entgegengesetzte Richtungen.
Dadurch bleibt der Gesamtimpuls erhalten.

\begin{PhysicsBox}[Materie wird zu Strahlung]
	\label{box:materie_wird_zu_strahlung}
	
	Bei der Annihilation werden Elektron und Positron nicht zu nichts.
	
	Ihre gesamte Energie bleibt erhalten und erscheint beispielsweise als
	Energie von Photonen:
	
	\begin{equation}
		e^-
		+
		e^+
		\longrightarrow
		2\gamma.
	\end{equation}
	
	Für ein nahezu ruhendes Paar besitzen die beiden Photonen jeweils eine
	Energie von ungefähr
	
	\begin{equation}
		\SI{511}{\kilo\electronvolt}.
	\end{equation}
\end{PhysicsBox}

\subsubsection*{Erhaltung der elektrischen Ladung}

Vor einer Paarbildung besitzt das Photon keine elektrische Ladung:

\begin{equation}
	q_{\gamma}
	=
	0.
\end{equation}

Nach der Paarbildung gilt

\begin{equation}
	q_{e^-}
	+
	q_{e^+}
	=
	-e
	+
	e
	=
	0.
\end{equation}

Die gesamte elektrische Ladung bleibt somit erhalten.

Auch bei der Annihilation gilt

\begin{equation}
	-e
	+
	e
	=
	0.
\end{equation}

Die entstehenden Photonen sind ebenfalls elektrisch neutral.

\begin{DidacticBox}[Antiteilchen und Erhaltungssätze]
	\label{box:antiteilchen_erhaltungssaetze}
	
	Teilchen und Antiteilchen besitzen entgegengesetzte elektrische
	Ladungen.
	
	Bei Paarbildung und Annihilation bleiben deshalb
	
	\begin{itemize}
		\item Energie,
		\item Impuls,
		\item Drehimpuls und
		\item elektrische Ladung
	\end{itemize}
	
	erhalten.
	
	Antiteilchen widersprechen den Erhaltungssätzen nicht. Sie sind
	vielmehr notwendig, damit diese Gesetze auch bei Teilchenerzeugung und
	Teilchenvernichtung erfüllt bleiben.
\end{DidacticBox}

\subsubsection*{Positronium}

Ein Elektron und ein Positron können vor ihrer Annihilation kurzzeitig
einen gebundenen Zustand bilden.

Dieser Zustand wird als Positronium%
\index{Positronium} bezeichnet.

Positronium ähnelt in seiner Struktur einem Wasserstoffatom. Anstelle
eines Protons und eines Elektrons besteht es jedoch aus

\begin{itemize}
	\item einem Elektron und
	\item einem Positron.
\end{itemize}

Da beide Teilchen dieselbe Masse besitzen, bewegen sie sich um ihren
gemeinsamen Schwerpunkt.

Positronium ist nicht dauerhaft stabil. Nach kurzer Zeit annihilieren
Elektron und Positron und erzeugen Photonen.

\begin{NoteBox}[Ein kurzlebiges Atom]
	\label{box:positronium_kurzlebiges_atom}
	
	Positronium ist ein gebundener Zustand aus Elektron und Positron.
	
	Es enthält keinen Atomkern. Beide Teilchen bewegen sich um den
	gemeinsamen Schwerpunkt und vernichten sich nach kurzer Zeit durch
	Annihilation.
\end{NoteBox}

\subsubsection*{Technische und medizinische Anwendung}

Die Annihilation von Elektronen und Positronen wird in der Medizin
genutzt.

Bei der Positronen-Emissions-Tomographie%
\index{Positronen-Emissions-Tomographie}
werden radioaktive Stoffe eingesetzt, die Positronen abgeben.

Ein ausgesandtes Positron verliert im Gewebe zunächst Bewegungsenergie.
Anschließend trifft es auf ein Elektron und annihiliert.

Dabei entstehen gewöhnlich zwei Photonen mit jeweils ungefähr

\begin{equation}
	E_{\gamma}
	=
	\SI{511}{\kilo\electronvolt}.
\end{equation}

Die Photonen verlassen den Entstehungsort in nahezu entgegengesetzten
Richtungen. Werden beide gleichzeitig registriert, kann die Lage des
Annihilationsereignisses bestimmt werden.

\begin{PhysicsBox}[Positronen in der medizinischen Bildgebung]
	\label{box:positronen_medizinische_bildgebung}
	
	Die Positronen-Emissions-Tomographie nutzt die Annihilation eines
	Positrons mit einem Elektron.
	
	Dabei entstehen zwei nahezu entgegengesetzt ausgesandte Photonen mit
	jeweils
	
	\begin{equation}
		\SI{511}{\kilo\electronvolt}.
	\end{equation}
	
	Durch die gleichzeitige Messung beider Photonen kann die räumliche
	Verteilung eines radioaktiven Markierungsstoffes im Körper bestimmt
	werden.
\end{PhysicsBox}

\subsubsection*{Bedeutung für das Verständnis des Elektrons}

Die negativen Energielösungen der Dirac-Gleichung führten zunächst zu
einem schwerwiegenden Interpretationsproblem.

Gerade dieses Problem offenbarte jedoch eine neue Eigenschaft der
Natur: Zu einem Teilchen kann ein Antiteilchen existieren.

Für das Elektron ist dieses Antiteilchen das Positron.

\begin{HistoryBox}[Von einer mathematischen Schwierigkeit zur Entdeckung]
	\label{box:mathematische_schwierigkeit_entdeckung_positron}
	
	Die negativen Lösungen der Dirac-Gleichung wirkten zunächst wie ein
	Fehler oder eine unerwünschte Folge der Theorie.
	
	Dirac nahm sie jedoch ernst.
	
	Aus ihrer physikalischen Deutung entstand die Vorhersage des Positrons,
	das wenige Jahre später experimentell nachgewiesen wurde.
	
	Die Geschichte zeigt, dass mathematisch notwendige Lösungen auf reale,
	bis dahin unbekannte Naturerscheinungen hinweisen können.
\end{HistoryBox}

Die Dirac-Gleichung machte damit deutlich, dass eine relativistische
Quantentheorie nicht dauerhaft auf eine feste Anzahl von Teilchen
beschränkt bleiben kann.

Elektronen und Positronen können erzeugt und vernichtet werden. Eine
vollständige Beschreibung solcher Prozesse erfordert deshalb den
Übergang von der Einteilchentheorie zur Quantenfeldtheorie.

